\chapter{Conclusion}

This thesis presented a reinforcement learning-based routing framework for improving the reliability of NoC interconnects under aging effects. The proposed approach integrates device-level aging models for EM and HCI into the routing decision process, enabling routers to dynamically adapt traffic distribution based on both performance and reliability considerations. Unlike prior aging-aware routing techniques that rely on fixed heuristics, the proposed method employs linear function approximation-based reinforcement learning to learn routing policies online during execution. This allows the routing strategy to adapt to changing traffic patterns and evolving network conditions without requiring explicit modeling of all possible scenarios.

Evaluation results demonstrate that the RL-based routing approach significantly improves system lifetime compared to deterministic routing, achieving up to $\sim5\times$ improvement in EM lifetime and consistent gains in HCI lifetime across multiple traffic patterns. When compared to the state-of-the-art Clotho-GAR approach, the proposed method achieves comparable or superior reliability improvements in most traffic patterns, particularly under asymmetric and structured traffic conditions.

In terms of performance, the RL-based approach maintains latency comparable to Clotho and consistently improves over DOR by delaying network saturation. While Clotho exhibits slightly better latency under certain traffic patterns such as transpose and bit reverse, the RL-based method provides more consistent performance across a wider range of traffic conditions. These results highlight a key trade-off. Heuristic-based routing can be optimized for specific patterns, whereas learning-based routing offers greater robustness and adaptability.

Overall, this work demonstrates that reinforcement learning is an effective and flexible framework for aging-aware routing in NoCs, offering improved reliability while maintaining competitive performance.